% Two drafts, both from the M3 measurement:
%   (1) replacement text for Section IV-F (what is actually quantized)
%   (2) a new Section VI subsection presenting Table~\ref{tab:m3:quant}, to be
%       placed after the deployment subsection; it is what the Section V audit
%       forward-link resolves to.
% Numbers: results/tables/tab_m3_quant.{csv,tex}, m3_quant_per_seed.csv
% (9 clean datasets; CSE-CIC-IDS2018 and Edge-IIoTset pending the repaired re-runs).
% Generated 2026-09-18. Labels to verify: sec:deployment, tab:m3:quant, subsec:audit.

% ===================== (1) Section IV-F replacement =========================
\paragraph{Deployed numerical format.}
The deployed student carries \emph{int8-grid weights executed in floating point}.
After training, every \texttt{Linear} weight tensor is projected onto a symmetric
per-tensor grid, $W_q = s\,\mathrm{clip}(\mathrm{round}(W/s), -127, 127)$ with
$s = \max|W|/127$, once and post hoc; no quantization-aware training, no
fine-tuning and no calibration set are used. Biases, activations and the
accumulator are untouched, and the projected values are stored and executed as
float32, so the grid constrains the weights without making inference integer
arithmetic. We therefore do not claim integer-only inference, and the energy
model of Section~\ref{sec:energy} prices synaptic operations, not int8 kernels.
What the projection buys is the storage figure quoted in
Section~\ref{sec:deployment} (one byte per weight when serialised) and the
guarantee that the reported accuracy is the accuracy of the weights that would be
shipped. Its cost is measured in Section~\ref{subsec:quantcost}.

% ===================== (2) new Section VI subsection ========================
\subsection{What Quantization Costs}
\label{subsec:quantcost}

Reporting a quantized artifact obliges us to say what the projection costs.
Table~\ref{tab:m3:quant} gives fp32 and int8-grid test macro-$F_1$ for all four
models on every benchmark, as 5-seed means, together with the mean and
worst-seed difference. For the two tree-guided models the floating-point side is
recovered by re-running training with the same seeds and capturing the weights
immediately before the projection; each such re-run reproduces the published
int8 number exactly, so the two columns differ only in the projection. Vanilla
SNN and TASNN are reported in floating point throughout the paper, so their int8
columns are counterfactual: what deployment would have cost had we quantized
them. For those two models we also apply the projection to every weight matrix
rather than to \texttt{Linear} layers alone, because Vanilla SNN keeps 98\% of
its weights in a recurrent cell that the \texttt{Linear}-only projection never
touches; for the tree-guided students, whose weights are all linear, the two
variants coincide.

Across the nine benchmarks unaffected by the dataset repairs of
Section~\ref{subsec:audit}, the projection is close to free for the models we
deploy: Tree-to-Spike averages $+0.06$ macro-$F_1$ points (worst dataset
$-0.42$, worst single seed $-1.82$) and DNN-Tree $-0.12$ (worst dataset $-0.38$,
worst seed $-2.22$). No dataset loses a full point in the mean for either model,
and the paired per-dataset difference between them is not significant
($p = 0.055$, Wilcoxon signed-rank, $N = 9$): on this evidence the spiking
student is never worse on average, but a claim that it is more robust than its
dense twin is not supported.
% PROVISIONAL -- do not freeze. p = 0.055 is over the nine clean datasets only.
% The repaired CSE-CIC-IDS2018 and Edge-IIoTset cells add two paired
% observations (N = 11) and can move p in either direction; the pre-repair
% CSE-CIC-IDS2018 DNN-Tree delta was -11.46, the largest in the table. Re-run
% the Wilcoxon from results/tables/m3_quant_per_seed.csv once
% scripts/run_m3_repaired.sh finishes, and reword this sentence to match: keep
% "no detectable difference" only if it still holds at N = 11.

That the cost is small is a property of these models, not of the method. The same
projection applied to the Vanilla SNN baseline costs $-0.47$ points on average,
$-1.82$ on UNSW-NB15 and as much as $-5.10$ on a single KDD Cup 99 seed, with two
of nine datasets losing more than a point. Per-tensor symmetric quantization is
safe here because the weight distributions of the distilled students are
well-conditioned -- their per-layer $\max|W| / p_{99}|W|$ ratios stay between 1.1
and 3.6 -- and not because post-training quantization is harmless in general. A
single outlier weight sets the scale for its whole tensor, and where one exists
the projection destroys the resolution of every other weight in that layer;
Section~\ref{subsec:audit} traces exactly such a case to a dataset artifact.

[RERUN] CSE-CIC-IDS2018 and Edge-IIoTset are re-run on the repaired corpora; their
rows in Table~\ref{tab:m3:quant}, and the accompanying statement of whether the
first-layer outlier survives the repair, are filled from those runs.
